\section{Water ambient and analytic perturbations}

The water extension keeps $z$ increasing into the incident medium and
$n=1-\ii\beta-\delta$.  The local glancing angle $\alpha_w$ specifies
\begin{equation}
 k_\parallel=k_0 n_w\cos\alpha_w,\qquad
 \kappa_j^2=k_0^2 n_j^2-k_\parallel^2,\qquad
 k_{z,j}=-\kappa_j.
\end{equation}
The slab and continuous solvers use the same radiation branch.  Incident
amplitude is normalized at a finite plane in homogeneous water.  Translating
this phase origin changes the complex amplitudes, including attenuation in an
absorbing ambient.  A measured external angle requires entrance-interface
propagation; it is not substituted for $\alpha_w$.

The local Green-function conversion becomes
\begin{equation}
 r_{\bm h}=\frac{2\pi\ii r_e}{A_{\rm ref} k_0 n_w\sin\alpha_f}
 F_{\bm h}.
\end{equation}
The normal components of the reciprocal p-polarized waves use
$k_0 n_w\cos\alpha/k_{z,j}$ consistently.  The squared amplitude ratio is a
local observable; macroscopic water-path attenuation and cell transmission
are separate.  In an absorbing incident half-space it must not be interpreted
as a globally conserved reflected-flux fraction.  See also S.~J.~Byrnes,
\emph{Multilayer optical calculations}, arXiv:1603.02720, for the distinction
between amplitude ratios and power flow with absorbing exterior media.

Water is split before integration:
\begin{equation}
 \rho_w=\rho_{w,0}+\Delta\rho_w,\qquad
 \Delta\rho_w\longrightarrow0\quad (z\longrightarrow+\infty).
\end{equation}
The equimolar sharp reference satisfies
$\int(\rho_w-\rho_\infty\Theta(z-z_G))\,dz=0$.
For the progressively broadened Gaussian stack,
$z_G=z_0+d(1/2-A_0)-\sigma_b^2/(2d)$.
Terrace averages translate both densities by the same heights and weights.
Only the finite, localized contrast enters the water matrix element; the
diagnostic atomic/reference arrays use a common finite baseline rather than
two separately divergent water half-space integrals.


The production reference contains only the water onset.  Sharp interfaces use
Fresnel coefficients; tanh interfaces use the exact isolated s-polarized
Epstein coefficients.  Neither adsorption peaks nor molecular layers introduce
additional optical segments.  The existing dry solid reference is retained,
adding one interface per physical terrace onset.

For molecular number density $\rho_m$ and $f(q)=f_O(q)+2f_H(q)$, the physical
transforms per lateral area are
\begin{align}
 W_{\rm step}(q)&=f(q)e^{\ii qz_0-\sigma_0^2q^2/2}
    \left(\frac{\ii\rho_m}{q}+\Gamma e^{\ii q\Delta z_a}\right),\\
 W_{\rm layered}(q)&=\rho_m d f(q)e^{\ii qz_0-\sigma_0^2q^2/2}
    \left[A_0-1+\frac{1}{1-e^{\ii qd-\sigma_b^2q^2/2}}\right].
\end{align}
The second expression sums the infinite damped Hosemann paracrystal.  Its
evaluation cost is independent of the number of visible water layers or their
damping length. Step adsorption has $\Gamma=A_0/A_{\rm cell}$, whereas layered
$A_0$ is the relative first-layer weight. The equimolar step position is
$z_G=z_0-\Gamma/\rho_m$ for the step family.

The reference transform and localized perturbation are
\begin{equation}
 W_0(q)=\frac{\ii\rho_m f(0)}{q}e^{\ii qz_G}B(q),\qquad
 D(q)=W(q)-W_0(q),\qquad
 B(q)=1\quad\hbox{or}\quad
 \frac{\sqrt{3}\sigma_0 q}{\sinh(\sqrt{3}\sigma_0 q)}.
\end{equation}
The implementation subtracts these half-spaces analytically using
\texttt{expm1} and small-$q$ series. It does not subtract two divergent numerical
quantities. The equimolar limit is $D(0)=0$; a frozen reference displaced from
the live equimolar plane instead gives the appropriate finite integrated excess.

For each terrace and each of four reciprocal plane-wave channels, $D(q)$ is
evaluated once at $q=-(\mu k_{z,i}+\nu k_{z,f})$. The coefficients include the
two local branch amplitudes, polarization contraction, and phase references.
The constant wavevectors on the water side of the frozen onset are extended
through the localized perturbation, as for an atomic form factor assigned to
its host medium. This is the declared local constant-medium DWBA approximation.
The finite coherent terrace sum remains; the infinite molecular sum is closed.
Off-specular direct water scattering vanishes because this water is laterally
uniform. The solid continues to use the native plane-wave DWBA kernel.

An isolated Epstein transition has
\begin{equation}
 \varepsilon(z)=\varepsilon_-+
 (\varepsilon_+-\varepsilon_-)\frac{1+\tanh((z-z_G)/w)}{2},
 \qquad w=\sqrt{12}\sigma_0/\pi.
\end{equation}
Writing $a=-\ii w k_{z,-}/2$ and $b=-\ii w k_{z,+}/2$, its center-referenced
downward-incidence s amplitudes can be written without critical-point gamma poles:
\begin{align}
 r&=\frac{b-a}{a+b}\frac{\Gamma(1-2b)}{\Gamma(1+2b)}
       \left[\frac{\Gamma(1+a+b)}{\Gamma(1+a-b)}\right]^2,\\
 t&=\frac{2b}{a+b}
       \frac{\Gamma(1+a+b)^2}{\Gamma(1+2a)\Gamma(1+2b)}.
\end{align}
Reverse-incidence coefficients interchange the upper and lower wavevectors.
Scattering-matrix composition connects these interface coefficients through
constant media. It does not sample hypergeometric internal fields or fit finite
tanh segments to a sampled density. This effective-interface construction is
not an exact Maxwell solution of overlapping tanh transitions. The plotted
onset envelope retains the first-order density-to-index convention and is not
a claim of an exactly solved global continuous profile.

Analytic tanh supports s polarization only. The production path explicitly
rejects p instead of silently invoking numerical integration; sharp interfaces
support both. For step-model confirmation, the explicitly selected erf path
solves the full physical step/adsorption profile numerically and evaluates
continuous DWBA matrix elements. It therefore tests the effective-interface and
constant-medium approximations as well as the reference choice, while retaining
identical physical densities and molecular factors. Its Maxwell equations are
\begin{equation}
 E''+(k_0^2\varepsilon-k_\parallel^2)E=0,\qquad
 (H'/\varepsilon)' +(k_0^2-k_\parallel^2/\varepsilon)H=0.
\end{equation}
Reciprocal electric fields enter bilinearly without complex conjugation.
Numerical erf convergence does not make the full atomistic DWBA exact.
